\begin{table}[htbp]
\centering
\small
\caption{Rebote tras barridos (brecha = O/E del lado contrario $-$ O/E base, $\delta$ = 2 bps; fuera de muestra)}
\label{tab:T4_rebote}
\begin{tabular}{lrrrr}
\toprule
Modelo & Vendedor, 1 h al alza & Vendedor, 1 h a la baja & Comprador, 1 h al alza & Comprador, 1 h a la baja \\
\midrule
E3 tesis (49 cov., f($\delta$)$\cdot$k(X)) & +0.223 & +0.054 & +0.027 & -0.157 \\
E3 f($\delta$)$\cdot$k(X), 44 cov. & +0.349 & +0.097 & -0.018 & -0.195 \\
E3 base ($-$k$\delta$, 44 cov.) & +0.469 & +0.195 & +0.164 & -0.019 \\
AS base & +1.799 & +1.024 & +2.427 & +0.608 \\
AS + hora + f($\delta$) & +1.749 & +1.083 & +2.688 & +0.608 \\
AS + 7 señales & -0.002 & +0.122 & +0.390 & -0.405 \\
GLM lineal 49 cov. & +0.276 & +0.216 & +0.307 & -0.092 \\
Hawkes 1 escala + hora + f($\delta$) & +1.745 & +1.077 & +2.676 & +0.606 \\
Hawkes 1 escala + señales & -0.014 & +0.119 & +0.389 & -0.404 \\
Hawkes 2 escalas + hora + f($\delta$) & +0.918 & +0.400 & +1.005 & +0.016 \\
Hawkes 2 escalas + señales & -0.187 & -0.112 & +0.016 & -0.516 \\
\bottomrule
\end{tabular}
\par\smallskip\footnotesize Positivo: llegan más agresores del lado contrario de los que el modelo espera; 0 = acierta como en momentos normales.
\end{table}
